% Generated draft: requires mathematical and source review.
\documentclass[11pt]{article}
\usepackage[margin=25mm]{geometry}
\usepackage{fontspec}
\setmainfont{Latin Modern Roman}
\usepackage{amsmath,amssymb,mathtools,unicode-math}
\setmathfont{Latin Modern Math}
\usepackage{enumitem,needspace,xurl,hyperref}
\setlength{\parindent}{0pt}
\setlength{\parskip}{5pt}
\setlength{\emergencystretch}{3em}
\newcommand{\N}{\mathbb{N}}
\newcommand{\Z}{\mathbb{Z}}
\newcommand{\Q}{\mathbb{Q}}
\newcommand{\R}{\mathbb{R}}
\newcommand{\C}{\mathbb{C}}
\newcommand{\F}{\mathbb{F}}
\newcommand{\A}{\mathbb{A}}
\newcommand{\PP}{\mathbb{P}}
\newcommand{\E}{\mathbb{E}}
\newcommand{\eps}{\varepsilon}
\newcommand{\dd}{\,\mathrm{d}}
\begin{document}
\section*{UnsolvedMath ID: 801044}
\textbf{Source problem code:} catalogue entry
\subsection{Definitions and mathematical statement}
Let \(\lambda\) denote Lebesgue probability measure on \([0,1]\). Under the usual measure-theoretic convention, \(G=\operatorname{Aut}([0,1],\lambda)\) is the group of invertible measure-preserving transformations modulo equality almost everywhere; inverses and compositions are understood modulo null sets. Define
\[
R=\{T\in G:\exists S\in G\ (T=S\circ S)\}.
\]
Thus \(T\in R\) means that some invertible measure-preserving transformation \(S\) satisfies \(T(x)=S(S(x))\) for \(\lambda\)-almost every \(x\in[0,1]\).

The source asks for a characterization of the transformations \(T\in R\) and states the conjecture
\[
R\text{ is }\boldsymbol{\Sigma}^{1}_{1}\text{-complete}.
\]
Here \(\boldsymbol{\Sigma}^{1}_{1}\) denotes the class of analytic sets. The source record does not specify the topology on \(G\), the reducibility used for completeness, or explicitly whether transformations are identified modulo null sets.

One conventional precise interpretation equips \(G\) with its weak Polish topology, characterized by
\[
T_n\longrightarrow T
\quad\Longleftrightarrow\quad
\lambda\bigl(T_n(A)\mathbin{\triangle}T(A)\bigr)\longrightarrow0
\quad\text{for every Lebesgue-measurable }A\subseteq[0,1],
\]
where \(\triangle\) denotes symmetric difference. Under continuous reductions from Baire space, the conjecture would say that \(R\) is analytic and that, for every analytic set \(B\subseteq\N^{\N}\), there exists a continuous map \(f:\N^{\N}\to G\) such that
\[
\forall x\in\N^{\N}\qquad
x\in B\ \Longleftrightarrow\ f(x)\in R.
\]
These topology, reduction, and null-set conventions are explicit interpretive choices requiring confirmation against the intended source formulation; they are not specified in the supplied record.
\subsection{Short English statement}
Determine which invertible Lebesgue-measure-preserving transformations of \([0,1]\) admit an invertible measure-preserving square root. The supplied source conjectures that this collection is \(\boldsymbol{\Sigma}^{1}_{1}\)-complete. Its precise descriptive-set-theoretic meaning requires confirmation of the topology and reduction convention. The registry labels the problem open; that label is supplied metadata, not an independently verified assessment.
\subsection{Sources}
\begin{itemize}
\item Local UnsolvedMath registry, permanent ID \(801044\), source problem code \texttt{catalogue entry}, title \emph{Open Problems in Mathematical Logic --- Problem 56 (Halmos, 1956)}. Supplied status: open. Upstream revision: \texttt{d371ee917ac9647fbb06cd6fc529deec6f3dda82}.
\item Barmpalias, Gao et al., \emph{Open Problems in Mathematical Logic} (2026), Problem 56. Supplied source URL: \url{https://arxiv.org/abs/2608.26628}. The record reports access on 2026-10-16 and extraction by \texttt{arxiv\_latex}; these are supplied provenance details, not a claim of independent access or verification.
\item The supplied attribution/context is ``Halmos, 1956.'' No full bibliographic reference for that attribution is supplied. The registry separately lists Barmpalias, Gao et al. as proposers and 2026 as the proposed year; these entries are retained as compilation metadata without resolving the historical attribution.
\end{itemize}
\end{document}
